\documentclass[a4paper,10pt]{article} \usepackage{anysize}
\marginsize{2cm}{2cm}{1cm}{1cm}
%\textwidth 6.0in \textheight = 664pt
\usepackage{xltxtra}
\usepackage{graphicx}
\usepackage{color} 
\usepackage{xgreek} 
\usepackage{fancyvrb}
\usepackage{minted}
\usepackage{listings}
\usepackage{enumitem} 
\usepackage{framed} 
\usepackage{relsize}
\usepackage{float} 
\setmainfont[Mapping=TeX-text]{DejaVu Serif}

%\setmainfont{Arial}
%\setsansfont{FreeSans}
%\setmonofont{FreeMono}
\begin{document}
%\def\thesubsection {\alph{subsection}}
\renewcommand{\labelenumi}{\alph{enumi}}
\renewcommand{\labelenumii}{(\roman{enumii})}

\begin{titlepage}
\begin{center}
\begin{figure}[t] 
     \includegraphics[scale=0.7]{title/ntua_logo}
\end{figure}
\begin{LARGE}\textbf{ΕΘΝΙΚΟ ΜΕΤΣΟΒΙΟ ΠΟΛΥΤΕΧΝΕΙΟ\\}\end{LARGE}
\vspace{5cm}
\begin{Large}
ΣΧΟΛΗ ΗΜ\&ΜΥ\\
Αλγόριθμοι και Πολυπλοκότητα\\
1\textsuperscript{η} Σειρά Γραπτών Ασκήσεων\\
Ακ. έτος 2010-2011\\
\end{Large}
\vspace{10cm}
\begin{tabular}{l r}
\Large{Λύρας Γρηγόρης}&
\large{Α.Μ.: 03109687}\\
\end{tabular}\\
\vspace{1cm}

\vfill
\large\today\\
\end{center}
\end{titlepage}



\section{Ασυμπτωτικός συμβολισμός, Αναδρομικές
Σχέσεις} \setcounter{section}{1}
\begin{enumerate}
\item Ταξινόμηση
\begin{enumerate}
\item $\log {n^3}$
\item $\sqrt{n}*\log^{50}{n} $
\item $\frac{n}{\log{\log{n}}} $
\item $\log{n!} $
\item $n*\log^{10}{n} $
\item $n^{1.01} $
\item $5*log_{2}{n} $
\item $\sum_{k=1}^n{k^5} $
\item $\log^{\log{n}}{n} = n^{\log{\log{n}}} $
\item $2^{\log_2^4{n}} $
\item $\log^{\sqrt{n}}{n} $
\item $e^{\frac{n}{\ln{n}}} $
\item $n*3^n $
\item $2^{2*n} $
\item $\sqrt{n!} $
\end{enumerate}

\item Τάξη Μεγέθους
\begin{enumerate}
\item $T(n)=5*T(n/7)+n*\log(n)$
$\Rightarrow n^{\log _7 {5}} = n^{0.827}$
$\Rightarrow n^{0.827} < n*\log{n}$\\
$\Rightarrow T(n) \in \Theta(n*\log{n})$

\item $T(n) = 4*T(n/5)+n/\log^2{n}$
$\Rightarrow n^{\log_5 {4}} = n^{0.861}$
$\Rightarrow n^{0.861} < n/\log^2{n}$\\
$\Rightarrow T(n) \in \Theta(n/\log^2{n})$

\item $T(n) = T(n/3) +3*T(n/7)+n$\\
$\Rightarrow T(n) \in O(n)$

\item $T(n) = 6*T(n/6)+n$
$\Rightarrow n^{\log _6 {6}} = n$\\
$\Rightarrow T(n) \in \Theta(n*\log {n})$

\item $T(n) = T(n/3) + T(2n/3) + n$\\
$\Rightarrow T(n) \in \Theta(n*\log{n})$

\item $T(n) = 16*T(n/4) + n^3*\log^2{n}$
$\Rightarrow n^{\log _4{16}} = n^2$
$\Rightarrow n^2 < n^3*\log^2{n}$\\
$\Rightarrow T(n) \in \Theta(n^3*\log^2{n})$

\item $T(n) = T(\sqrt(n)) + \Theta(\log{\log{n}})$
$(k=\log{n},f(k)=T(n))$\\
$\Rightarrow
T(\sqrt{n})=f(\log{\sqrt{n}})=f(\frac{1}{2}*\log{n})=f(\frac{k}{2}) $\\
$\Rightarrow f(k)=f(\frac{k}{2})+\Theta(\log{k})$\\
$\Rightarrow \Theta(n)$

\item $T(n) = T(n-3) + \log{n}$\\
$\Rightarrow T(n) \in \Theta(n*logn)$

\end{enumerate}
\end{enumerate}

\section{Ταξινόμιση σε Πίνακα με Πολλά Ίδια Στοιχεία}

\section{Δυαδική Αναζήτηση}
\begin{enumerate}
\item Ξεκινάμε με ένα στοιχείο. Και συγκρίνουμε με τον $x$. Όσο ο $x$ είναι
μεγαλύτερος από αυτόν διπλασιάζουμε τον αριθμό των στοιχείων και ελέγχουμε
πάλι με τον τελευταίο. Μόλις φτάσουμε σε μεγαλύτερο αριθμό από τον $x$ έστω
στη θέση $k$ και εφαρμόζουμε κλασσική δυαδική αναζήτηση στο τμήμα $A[1...k]$. 
Μέγιστος αριθμός επαναλήψεων είναι $O(\log{k})$.
\item Παίρνω τα k πρώτα στοιχεία από κάθε πίνακα. Τα χωρίζω στη μέση και
συγκρίνω τα $A[k/2]$ με $B[k/2+1]$ και $A[k/2+1]$ με $B[k/2]$. Αν
$A[k/2]<B[k/2+1]$ και $B[k/2]<A[k/2+1]$ τοτε το k-οστό στοιχείο είναι το
$max\{A[k/2],B[k/2]\}$. Αν $A[k/2]<B[k/2+1]$ και $B[k/2]>A[k/2+1]$ τότε
επαναλαμβάνω χρησιμοποιώντας τα στοιχεία $[A[k/2+1]..A[k]]$ και
$[B[1]..B[k/2]]$.
\end{enumerate}
\section{Συλλογή Comics}
Ζητάμε πρώτα το MSB για όλα τα τεύχη. Έτσι τα χωρίζουμε σε $\frac{n}{2}$ και
$\frac{n}{2}-1$. Προφανώς αυτό που λείπει είναι στα λιγότερα. Εφαρμόζουμε πάλι
κάνοντας $\frac{n}{2}$ ερωτήσεις και κρατάμε πάλι το υποσύνολο με το μικρότερο
πλήθος. Αναδρομικά θα καταλήξουμε στο τεύχος που λείπει από τη συλλογή έχοντας
ρωτήσει $n+\frac{n}{2}+\frac{n}{4}+..+1=2*n*(1-\frac{1}{2^n}) < 2*n$ φορές.
\section{Πολυκατοικίες χωρίς Θέα}
Ξεκινάμε με τον από τον $A[1]$ θέτοντας το $B[1]=0$ εφόσον δεν υπάρχει κάποιος
δυτικότερα από αυτόν. Προχωράμε στον επόμενο $A[2]$ και συγκρίνουμε το ύψος
του με τον $A[1]$. Αν $A[2]<A[1]$ τότε θέτουμε $B[2]=1$ αλλιώς συγκρίνουμε το
ύψος $A[2]$ με το ύψος του στόχου του $A[1]$. Αν πάλι δεν βρούμε ψηλότερο
τρέχουμε πάλι για το στόχο του στόχου αναδρομικά μέχρι να φτάσουμε σε 0
ή σε κάποιον ψηλότερο του $A[2]$ και τον θέτουμε ως στόχο στο πεδίο $B[2]$.
Επαναλαμβάνουμε μέχρι $n$. Μέγιστος αριθμός συγκρίσεων είναι $2n$.

\end{document}
